\documentclass[12pt,a4paper]{report}
% \usepackage[utf8]{inputenc}   % Sin compilador Unicode
\usepackage[english]{babel}     % Poder utilizar caracteres españoles
\usepackage{graphicx}           % Importar imágenes
\usepackage{fancyhdr}           % Modificar la cabecera y el pie de página
\usepackage[a4paper, margin=1in]{geometry} % Márgenes más ajustados
\usepackage[dvipsnames]{xcolor} % Emplear colores
\usepackage{caption}            %
\usepackage{subcaption}         % Subfigures
\usepackage{amsmath,amssymb}    % Fórmulas matemáticas
\usepackage{tabu}               % Formateo de tablas
\usepackage{array}              % Más formateo de tablas
\usepackage{listings}           % Código
\usepackage{mdframed}           % Bordes bonitos
\usepackage[backend=biber, style=numeric]{biblatex} % Referencias
\usepackage{csquotes}
\usepackage{lastpage}           % Referenciar última página
\usepackage[toc, acronym]{glossaries}% Glosario
\usepackage[
	colorlinks=true,
	linkcolor=black,
	urlcolor=NavyBlue,
	citecolor=green,
	anchorcolor=black]{hyperref}    % Hipervínculos con colores


% COLORES ------------------------------------------------------------------------------------------
\definecolor{light-gray}{gray}{0.95} % Color de stack overflow
\definecolor{mygreen}{rgb}{0,0.6,0}
\definecolor{mymuave}{rgb}{0.58,0,0.82}

% BIBLIOGRAFÍA -------------------------------------------------------------------------------------
%\addbibresource{./referencias.bib}

% FONTS --------------------------------------------------------------------------------------------
\renewcommand{\familydefault}{\sfdefault}

% ESPACIADO ----------------------------------------------------------------------------------------
\setlength{\parskip}{1em}

% NUEVA PÁGINA POR SECCIÓN -------------------------------------------------------------------------
\let\oldsection\section
\renewcommand\section{\clearpage\oldsection}

% CABECERA Y PIE DE DOCUMENTO ----------------------------------------------------------------------
\fancypagestyle{myheader}{
	\fancyhf{} % Clear all headers/footers
	\fancyhead[R]{\MakeUppercase{\leftmark}} % Cabecero a la derecha
	\fancyhead[L]{\MakeUppercase{C++ Engine Architecture}}
	\fancyfoot[R]{\textcolor{gray}{PAGE}~\thepage\ |~\pageref{LastPage}} % Número de página
	\fancyfoot[L]{\includegraphics[width=1em]{imgs/gobbly-icon.png}} % Icono u-tad
	\renewcommand{\headrulewidth}{1pt} % 1pt header rule
	\renewcommand{\headrule}{\
		\hbox to\headwidth{\color{NavyBlue}\leaders\hrule height\headrulewidth\hfill}}
	\renewcommand{\footrulewidth}{0pt} % No footer rule
}

% FORMATO CÓDIGO -----------------------------------------------------------------------------------
\renewcommand{\lstlistingname}{Código}% Listing -> Código
\renewcommand{\lstlistlistingname}{Lista de \lstlistingname\ s}% Lista de Códigos
\lstset{
	backgroundcolor=\color{light-gray},
	basicstyle=\footnotesize,
	commentstyle=\color{mygreen},
	escapeinside={\%*},
	extendedchars=true,
	firstnumber=1,
	keywordstyle=\color{NavyBlue},
	numbers=left,
	numbersep=5pt,
	numberstyle=\tiny,
	rulecolor=\color{black},
	tabsize=2,
	showstringspaces=false
}

% GLOSARIO -----------------------------------------------------------------------------------------
% \makeglossaries
% \newglossaryentry{test}
% {
%		name={test},
%		descriptioin={This is a test},
%		plural={tests}
% }

% DOCUMENTO ========================================================================================
\begin{document}
\pagestyle{empty}

% PÁGINA TÍTULO ------------------------------------------------------------------------------------
\begin{titlepage}
	\begin{center}
		\vspace*{1cm}

		\Huge
		\textbf{C++ Engine Architecture}

		\vspace{0.5cm}
		\LARGE
		An Entity-Relation Model

		\vspace{1.5cm}

		\textbf{Itziar Morales Rodríguez}

		\vfill

		\vspace{0.8cm}

		\includegraphics[width=0.4\textwidth]{imgs/cpp-logo.png}

		\Large
		Software Engineering and Computational Mathematics\\
		U-TAD\\
		Spain\\
		\today

	\end{center}
\end{titlepage}
\pagebreak

% ABSTRACTO ----------------------------------------------------------------------------------------
\begin{abstract}
	% TODO
\end{abstract}
\pagebreak

% ÍNDICE -------------------------------------------------------------------------------------------
\tableofcontents

% SECCIONES ========================================================================================
\pagebreak
\pagestyle{myheader}
\setcounter{page}{1}

% INTRODUCCIÓN -------------------------------------------------------------------------------------
\section{Introduction} % TODO
\par

% GLOSARIO =========================================================================================
% \printglossaries

% BIBLIOGRAFÍA =====================================================================================
% \printbibliography[heading=bibintoc, title={Referencias}]

\end{document}
